%% ─────────────────────────────────────────────────────────────────────────
%%  ARR "Explanation of Revisions" — required at resubmission.
%%  Resubmission of ARR submission 1656.
%%  Compile:  pdflatex Revision_Notes.tex  (twice)
%% ─────────────────────────────────────────────────────────────────────────
\documentclass[11pt]{article}
\usepackage[margin=1in]{geometry}
\usepackage[T1]{fontenc}
\usepackage[utf8]{inputenc}
\usepackage{times}
\usepackage{amsmath}
\usepackage{microtype}
\usepackage{booktabs}
\usepackage{longtable}
\usepackage{array}
\usepackage{xcolor}
\usepackage[hidelinks]{hyperref}

\newcommand{\done}{\textcolor{blue}{\textbf{Addressed}}}
\newcommand{\newexp}{\textcolor{blue}{\textbf{New experiment}}}
\newcommand{\scoped}{\textcolor{blue}{\textbf{Scoped}}}

\setlength{\parindent}{0pt}
\setlength{\parskip}{5pt}

\title{\vspace{-2em}Explanation of Revisions\\
\large Resubmission of ARR submission 1656:\\
``Correction and Corruption Signals: A Controlled Study of
Capability-Asymmetric Multi-Agent LLM Debate\\ Under Adversarial Wrong-Peer
Exposure''}
\author{}
\date{}

\begin{document}
\maketitle
\vspace{-3em}

\section*{Summary}

We thank the three reviewers and the Area Chair. The meta-review's decision
text identified the crux precisely:

\begin{quote}
``Although the authors stated in their response that additional model runs
have been completed, in-cycle ACL policies reserve these for future revisions,
leaving the current paper scope under-supported.''
\end{quote}

Every experiment the reviewers asked for had been run but could not be
reported in an author response. This revision puts all of them
\emph{in the manuscript}, on the identical question pool, random seed, prompts,
and judge cascade as the reviewed version, and adds three further experiments
that go beyond what was requested.

\paragraph{What is new in the manuscript.}
\begin{enumerate}\itemsep2pt
  \item \textbf{A solo re-answering control (C1R)} --- the causal baseline
        sTG4 and the AC identified as the central methodological gap. Section~4.2.
  \item \textbf{An honest-peer control (C3H, C4H)} --- the same weak models
        answering naturally, addressing Tf6M's and sTG4's objection that
        persona-anchored peers are adversarial distractors rather than natural
        low-capability output. Section~4.2.
  \item \textbf{A confidence-filter re-run with confidence actually elicited
        (C5R, C5H)} plus a corrected pre-flight test --- the ``properly designed
        follow-up'' Tf6M asked for. Section~3.4, Section~4.4.
  \item \textbf{An eight-model capability sweep} --- the systematic sweep
        including open-weight models that Tf6M argued was necessary to show a
        general phenomenon rather than a pairwise anomaly. Section~4.3,
        Table~5, Figure~1.
  \item \textbf{GPT-4o-mini extended from 50 to the full 300-question pool},
        answering gL73's concern about limited questions.
  \item \textbf{A contamination probe} on numerically perturbed GSM8K.
  \item \textbf{A controlled split-peer condition (C4split)} --- tests whether
        a visible correct peer changes harmful revision, removing the last
        ``deferred to future work'' in the previous version. Section~4.6.
  \item \textbf{A heterogeneous-smart debate control (C2het)} --- three
        architecturally distinct strong models, converting the
        ``homogeneous debate is close to sampling'' limitation into a result.
  \item \textbf{An output-distribution probe} on an open-weight focal model,
        measuring the probability mass that moves onto the peer-asserted wrong
        answer between rounds --- evidence beneath the discrete choice. We call
        it an output-distribution probe, not a mechanism: it observes output
        probabilities and does not identify a representation-level cause.
\end{enumerate}

\paragraph{What is new in the writing.}
The term \emph{focal model} is now defined at first use in both the abstract
and Section~1 (gL73~W1); Section~1 is restructured around three explicit
research questions and the contributions each end with a pointer to the
section, table or figure that supports it
(gL73~W2, W5); the conditions are given in a table rather than in prose; and
Section~2.2 gains a paragraph on \textbf{debate as an error detector}, the
specific prior-work area the meta-review named as under-discussed (M7,
gL73~W4).

\paragraph{Corrections we made to the previous version.}
In preparing this revision we re-derived every number in the manuscript
directly from the released per-trial logs, and corrected fourteen
discrepancies that survived the reviewed version. The full list is in the
released \texttt{audit/numerical\_consistency\_report.md}; the ones that change
a reported number are:
\begin{itemize}\itemsep2pt
  \item Table~3 previously reported \emph{raw} McNemar $p$-values under a
        caption stating they were Holm-corrected. Both are now reported in
        separate columns. All conclusions that were claimed as significant
        remain significant after correction.
  \item The contamination-probe contrast was reported as $p=0.07$ without
        noting that this question-level test rests on eight discordant pairs.
        Section~4.5 now states plainly that the probe is underpowered at the
        question level ($p=0.23$ after correction) while significant at the
        trial level, and no longer claims significance for it.
  \item Appendix~S1's trial-count table still carried the superseded
        50-question GPT-4o-mini numbers, which contradicted Table~2. All data
        tables are now generated programmatically from the released logs by a
        released script, so they cannot drift from the data again.
  \item Appendix~S7 asserted that the weak models' unprompted accuracy was
        ``not measured in this study''. The honest-peer conditions measure it
        directly; the numbers are now given.
  \item The released corrected-gate report substituted zero for
        $P(\text{loud}\mid\text{correct})$ on the wrong-anchored substrate,
        where that class is empty by construction. This produced an apparent
        discriminative gap of $0.99$ and a spurious ``passed'' verdict in the
        artefact, contradicting the paper's own table. The report now records
        the substrate as undefined and gives per-class counts.
  \item The capability-sweep Spearman $\rho$ was computed on the 1-dp rounded
        table values, which tie four mid-range models together and inflated it
        to $-0.97$ with $p=0.0001$. On unrounded rates it is $\rho=-0.95$ with
        exact permutation $p=0.0011$.
  \item Persona validation counts in Appendix~S5 were each off by 14 (the
        failures were double-counted): the correct split is $1{,}116$ passing
        first, $370$ after regeneration, $14$ failing.
  \item The primary focal model was described as being of the DeepSeek-V3
        family with a citation to that technical report. The API resolves the
        requested alias to the served snapshot \texttt{deepseek-v4-flash}, for
        which no such report exists, so the manuscript now reports the served
        snapshot recorded in the run metadata and makes no architectural claim.
\end{itemize}

\clearpage
\section*{Point-by-point response}

\subsection*{Meta-review (Area Chair QTRJ)}

\begin{longtable}{@{}p{0.30\textwidth}p{0.14\textwidth}p{0.50\textwidth}@{}}
\toprule
\textbf{Point} & \textbf{Status} & \textbf{Where it is answered} \\
\midrule
\endhead
Only two focal models; a systematic sweep across a wider set, including
open-weight models, is necessary to show a general phenomenon.
& \newexp
& Eight focal models spanning $44.9$--$76.2\%$ solo accuracy, on the same
  300-question pool. All eight point estimates of the accuracy change under two
  confidently wrong peers are non-negative, and harmful revision rises steeply
  as solo accuracy falls (Spearman $\rho=-0.95$, exact permutation $p=0.001$;
  leave-one-model-out $\rho\in[-0.96,-0.93]$; $\rho=-0.94$ excluding the two
  weakest). Section~4.3, Table~5, Appendix Figure. \\
\addlinespace
Persona prompts turn weak peers into deliberate adversarial distractors;
the manuscript remained scoped to distractors.
& \newexp\ \scoped
& Honest-peer conditions (C3H, C4H) run the same weak models with the standard
  prompt. Their unprompted accuracy is $50.4\%$ and $45.1\%$. Honest peers leave
  focal accuracy at solo level ($-0.3$~pp, $p=1.00$) and are clearly below the
  wrong-anchored condition ($-7.3$~pp, $p=0.004$ Holm). The claims are scoped to
  confidently wrong peers throughout. Section~3.1, Section~4.2. \\
\addlinespace
Major flaw in causal attribution: the $+7.9$~pp gain cannot be attributed to
peer exposure without a re-examine/re-answer control.
& \newexp
& C1R gives the focal model the Round-1 revision instruction with no peer
  messages. Re-answering alone is indistinguishable from solo
  ($+0.3$~pp, $p=1.00$); adding two wrong peers on top of that same second
  attempt gains $+6.7$~pp ($p=0.011$ Holm). Section~4.2, Table~3. \\
\addlinespace
C5 removed 100\% of peer messages because the weak-peer prompt omitted the
Round-0 confidence instruction; this renders C5 uninformative.
& \newexp
& The persona prompt was repaired and the filter re-run on peers that do state
  confidence (C5R), plus an honest variant (C5H). Peers now emit confidence in
  essentially every Round-0 message. The filter still removes almost nothing,
  now for a \emph{signal} reason rather than a parsing one: confidence does not
  separate right from wrong (AUROC $0.58$). Section~3.4, Section~4.4,
  Appendix~S3, Appendix~S4. \\
\addlinespace
``Focal agent'' left undefined in the main text.
& \done
& Defined at first use in the abstract and in Section~1. \\
\addlinespace
Presentation of the research questions was unclear.
& \done
& Explicit RQ1--RQ3 in Section~1, each mapped to the section that answers it. \\
\addlinespace
Prior work on collaborative debate in error detection insufficiently
discussed.
& \done
& New paragraph in Section~2.2 covering cross-examination as a factual-error
  detector, divergence-preserving debate, and the self-correction result that
  motivates external signal, with an explicit contrast: those works treat the
  peer as an examiner whose disagreement flags an error, whereas we find a
  \emph{confidently wrong} peer supplies a usable prompt to recompute. \\
\bottomrule
\end{longtable}

\subsection*{Reviewer gL73}

\begin{longtable}{@{}p{0.30\textwidth}p{0.14\textwidth}p{0.50\textwidth}@{}}
\toprule
\textbf{Weakness} & \textbf{Status} & \textbf{Response} \\
\midrule
\endhead
W1. ``What is the `focal agent'?''
& \done
& The focal model is the single designated agent whose Round-0 and Round-1
  answers are scored; all other agents are peers whose Round-0 messages it sees.
  Defined in the abstract and at first use in Section~1. \\
\addlinespace
W2. The writing is poor and the main research problems are unclear.
& \done
& Section~1 rebuilt around RQ1--RQ3; the five-condition setting moved from prose
  into Table~1 with a plain-language ``what it answers'' column; statistics moved
  out of prose into tables; the abstract rewritten without condition codes. \\
\addlinespace
W3. Experiments limited to a limited number of questions and LLMs.
& \newexp
& Eight focal models instead of two, and GPT-4o-mini extended from 50 to all
  300 questions ($1{,}500$ trials per condition instead of $250$). The headline
  harmful-flip rise for GPT-4o-mini is now a well-powered $7.4\%\to13.4\%$. \\
\addlinespace
W4. A closely related work has not been sufficiently discussed.
& \done
& We asked which work was meant; the meta-review identified the area as
  collaborative debate for error detection. Section~2.2 now discusses it (see
  M7 above). Section~2.2 also already added adversarial and
  capability-asymmetric debate. \\
\addlinespace
W5. Unclear how the experiments support the listed contributions.
& \done
& Three contributions instead of four, each ending with the section, table or
  figure that supports it. \\
\addlinespace
Datasets $=1$, Software $=1$, Reproducibility $=2$.
& \done
& The full pipeline, all prompts, the question and persona pools, the per-trial
  logs for every trial, and the scripts that regenerate every number and figure
  are released now rather than on acceptance, through an anonymised repository
  linked in the Reproducibility section. The data tables in the paper are
  generated from those logs by a released script. \\
\bottomrule
\end{longtable}

\subsection*{Reviewer Tf6M}

\begin{longtable}{@{}p{0.30\textwidth}p{0.14\textwidth}p{0.50\textwidth}@{}}
\toprule
\textbf{Point} & \textbf{Status} & \textbf{Response} \\
\midrule
\endhead
W1. The main claim rests on only two focal models; a wider sweep including
open-weight models is needed to establish a general phenomenon rather than a
pairwise difference. The paper does not explore whether the difference is
predictable from a capability metric.
& \newexp
& The sweep spans eight focal models, six of them open-weight. The qualitative
  difference \emph{is} predictable from a capability metric: solo accuracy
  predicts the harmful-flip rate almost monotonically
  (Spearman $\rho=-0.95$, exact permutation $p=0.001$, $n=8$), from $6.5\%$ for the strongest
  model to $24.9\%$ for the weakest. Notably the \emph{gain} does not track
  ability ($\rho=0.33$, $p=0.42$) --- it is the cost, not the benefit, that is
  graded by strength, which is a sharper claim than the reviewed version made.
  Section~4.3, Table~5. \\
\addlinespace
W2. Anchored peers should be read as adversarial distractors, not honest
low-capability outputs.
& \scoped\ \newexp
& We agree, and now demonstrate the distinction rather than only stating it.
  The honest-peer control shows the correction effect is specific to
  confidently wrong peers. Every conclusion is scoped accordingly, in the
  abstract, Section~1, Section~4.2 and Limitations. \\
\addlinespace
W3. The confidence-weighting experiment is uninformative because the persona
prompt omits the confidence line, giving 100\% peer exclusion; the failure
cannot be attributed to the filtering mechanism.
& \newexp
& Correct, and now fixed. With confidence elicited, the filter is tested on
  peers that state it, and still fails --- because weak peers are equally loud
  when right and when wrong. We also give a corrected pre-flight test on the
  discriminative gap
  $P(\text{loud}\mid\text{wrong})-P(\text{loud}\mid\text{correct})$, which fails
  in advance and would have predicted the null, whereas the original
  loud-when-wrong test passes at $0.995$ and does not. \\
\addlinespace
Suggestions 1--4 (re-answering control; honest weak peers; re-run the filter
with reliable Round-0 confidence; expand to more focal models).
& \newexp
& All four carried out and reported: C1R; C3H/C4H; C5R/C5H with the corrected
  gate; the eight-model sweep. \\
\bottomrule
\end{longtable}

\subsection*{Reviewer sTG4}

\begin{longtable}{@{}p{0.30\textwidth}p{0.14\textwidth}p{0.50\textwidth}@{}}
\toprule
\textbf{Point} & \textbf{Status} & \textbf{Response} \\
\midrule
\endhead
W1. The study is not about natural capability asymmetry --- the weak peers are
deliberately wrong distractors, not natural mistakes.
& \scoped\ \newexp
& The honest-peer conditions supply the natural-mistake comparison directly,
  and the claims are re-scoped. The honest condition also yields a
  composition gradient that the anchored design cannot produce: harmful flips
  run $25.9\%$ when both peers happen to be wrong, $7.9\%$ when they disagree,
  and $1.4\%$ when both are right. \\
\addlinespace
W2 (main weakness). The $+7.9$~pp gain cannot be attributed to peer exposure
without a control where the focal model simply re-examines the question.
& \newexp
& This is exactly C1R, and it is now the paper's second experiment. Re-answering
  alone reproduces none of the gain. We had noted in our response that C5's
  total peer removal made an accidental partial control; that argument is
  withdrawn in favour of the dedicated, powered experiment the reviewer asked
  for. \\
\addlinespace
Suggestion 1: solo re-answering control with the same Round-1 instruction and
reasoning budget but no peer messages.
& \newexp
& Implemented exactly as specified: identical Round-1 instruction and token
  budget, peer block absent. Section~3.1, Section~4.2. \\
\addlinespace
Suggestion 2: evaluate honest weak peers not anchored to wrong answers.
& \newexp
& C3H and C4H, $3{,}000$ trials. \\
\addlinespace
Suggestion 3: rerun confidence filtering with a prompt that reliably elicits
Round-0 confidence from all peers.
& \newexp
& C5R and C5H. \\
\addlinespace
Suggestion 4: expand the analysis to more focal models with different
capability.
& \newexp
& Eight focal models. \\
\bottomrule
\end{longtable}

\section*{Beyond what was requested}

Three experiments in this revision were not asked for. They close the two
weaknesses the previous version conceded in its own Limitations, and add
mechanistic evidence.

\begin{itemize}\itemsep3pt
  \item \textbf{Split-peer control (C4split), 1{,}500 trials.} The reviewed
        version admitted its design could not separate ``a wrong answer is
        present'' from ``every visible peer is wrong'', and deferred a
        split-peer control to future work. It is now run: one wrong-anchored
        and one correct-anchored peer, same pool and replications. Adding a correct
        peer lowers harmful revision ($4.8\%\to3.4\%$ paired, $n=1{,}042$,
        $p=0.03$) \emph{and} raises accuracy to $86.3\%$ --- $+4.7$ points over
        two wrong peers ($p=0.013$ corrected) and $+11.7$ over solo --- while
        the correction rate is unchanged ($56.0\%$ vs $55.8\%$).
        We initially read this as evidence that \emph{unanimity} drives the
        harm, then tested that directly and found it false: within C4 the two
        anchored peers name the same wrong answer in only $16.4\%$ of trials,
        and whether they agree makes no difference to harmful revision
        ($5.95\%$ vs $6.58\%$, $p=0.88$). The manuscript therefore claims only
        that a \emph{visible correct peer} reduces harmful revision, and states
        that C4split changes two things at once (removes a wrong peer, adds a
        correct one) so it isolates neither.
  \item \textbf{Heterogeneous-smart debate (C2het), 900 trials.} The reviewed
        version conceded that its homogeneous condition was closer to
        self-consistency sampling than to debate (pairwise $\kappa=0.72$,
        ESS $\approx1.2$). Three architecturally distinct strong models now
        debate on the same pool. This heterogeneous control did not improve
        performance: accuracy is indistinguishable from solo ($+0.7$,
        $p=0.87$), the focal ends below where it started
        ($79.0\%\to77.3\%$) where self-sampling gains ($75.8\%\to78.7\%$), and
        harmful revision rises from $3.3\%$ to $5.8\%$ ($n=638$, $p=0.02$).
        Because both distinct peers are also $9$--$10$ points weaker than the
        focal, diversity and peer strength are not separated and the control
        does \emph{not} estimate the effect of diversity alone; the manuscript
        says so explicitly and keeps the corresponding limitation.
  \item \textbf{Mechanistic log-probability probe.} Limitations previously said
        the study was ``behavioural, not mechanistic''. On an open-weight focal
        model served locally (\textit{Llama-3.1-8B-Instruct}, $2{,}700$ trials)
        we measure the probability mass that moves onto the peer-asserted wrong
        answer between rounds. Homogeneous peers move the model \emph{away}
        from that answer ($-0.018$); two confidently wrong peers move it
        \emph{toward} it ($+0.064$). Paired by (question, replication) the
        shift is $+0.082$ (95\% CI $[+0.067,+0.097]$, Wilcoxon
        $p\approx10^{-42}$, $n=898$). The pull is therefore present and
        measurable on trials where the stated answer never changes --- which
        flip rates structurally cannot show.
\end{itemize}

\paragraph{One finding that cuts against our own design.}
The honest-peer control turned out to do more than establish scope. Comparing
trials where both weak peers were wrong, paired by question so difficulty is
held fixed, peers that are wrong \emph{naturally} flip the focal model about
twice as often as peers \emph{instructed} to be wrong ($29.7\%$ against
$14.4\%$, $p<10^{-5}$). This survives stratifying on whether the two peers
named the same wrong answer, so it is not an artefact of consensus strength.
The wrong-answer personas argue in exaggerated faulty styles that are
comparatively easy to discount. We now state plainly that our anchored
conditions are \emph{conservative} as a probe of corruption risk, and that a
real mixed-ability ensemble is likely to be more dangerous than the adversarial
one we built --- the opposite of the usual concern about adversarial framing,
and a direct answer to Tf6M~W2 and sTG4~W1.

\section*{Reviewer and meta-reviewer preference}

We request \textbf{the same reviewers and the same Area Chair}. Every point
raised is implemented in the manuscript rather than promised, on the identical
pool, seed and protocol, and the returning reviewers are best placed to judge
whether the prior weaknesses are resolved. The title is unchanged so the paper
is recognisable.

\end{document}
